\documentclass[12pt]{article}
\usepackage[margin=1in]{geometry}
\usepackage{setspace}
\usepackage{amsmath, amssymb, amsthm}
\usepackage{natbib}
\usepackage{booktabs}
\usepackage{hyperref}
\usepackage{enumitem}

\title{Political Polarization, Institutional Trust, and the Interpretation of Auditor Signals}
\author{Neil Hwang\\
        \small Bronx Community College, CUNY}
\date{\today}

\newtheorem{hyp}{Hypothesis}
\newtheorem{proposition}{Proposition}
\newtheorem{corollary}{Corollary}

\begin{document}
\maketitle
\doublespacing

\begin{abstract}
This paper studies whether political polarization changes how investors interpret
standardized auditor disclosures. We focus on auditor changes disclosed under
Item 4.01 of Form 8-K, a setting in which all investors observe the same
mandated public signal at the same time. We argue that polarization matters not
because it alters the disclosure itself, but because it increases heterogeneity
in investors' trust in auditors and in the regulatory institutions that make
those disclosures credible. In a Bayesian framework with heterogeneous priors,
polarization widens posterior-belief dispersion after a common audit signal,
thereby increasing disagreement-driven trading volume and the magnitude of price
responses. The effect is strongest when the disclosure is more ambiguous, so
that priors matter more relative to signal precision. Empirically, we test
whether market reactions to auditor-change disclosures are stronger in more
polarized political environments. The paper's central contribution is to connect
three literatures that have largely developed separately: audit disclosures in
capital markets, polarization in finance, and heterogeneous-belief asset
pricing. Rather than claiming that polarization affects financial disagreement in
general, we show how polarization changes the interpretation of auditor-originated
signals in a highly regulated disclosure setting. This sharper mechanism helps
position the paper as a contribution to auditing and capital markets research,
with direct implications for disclosure-based audit regulation.
\end{abstract}

\section{Introduction}
\label{sec:intro}

Auditor-related disclosures are intended to provide a standardized,
institution-backed signal about the reliability of firms' financial reporting.
Yet market reactions to those disclosures vary substantially across otherwise
similar events. This paper argues that one important source of that variation is
political polarization. Specifically, we propose that polarization increases the
cross-sectional dispersion of investor priors about the credibility of auditors,
audit oversight, and the institutions that support financial reporting.
Consequently, even when investors observe the same public audit signal at the
same time, they interpret that signal less uniformly in more polarized
environments.

We study auditor changes disclosed under Item 4.01 of Form 8-K. This setting is
well suited to our question for three reasons. First, the disclosure is
mandatory, highly standardized, and publicly released to all investors
simultaneously. That feature sharply limits explanations based on differential
information access. Second, auditor changes are often intrinsically ambiguous. A
change can reflect routine auditor rotation, fee considerations, or firm growth,
but it can also reflect disagreements over accounting issues, internal control
concerns, or deteriorating reporting quality. Investors therefore must interpret
the event rather than merely record it. Third, the economic meaning of an
auditor-change disclosure depends critically on institutional trust. Investors
must assess whether the outgoing and incoming auditors are credible monitors,
whether the firm is being transparent about the reasons for the change, and how
much weight to place on the surrounding regulatory regime. Those judgments are
especially likely to differ when political polarization erodes common beliefs
about institutional reliability.

The paper's core idea is narrower and more defensible than the broad claim that
``polarization affects financial markets.'' A growing finance literature already
shows that polarization affects expectations, disagreement, and trading in a
variety of settings \citep{kempf2024, bordalo2023, hao2026}. Our contribution is
to show how that logic operates in an auditing setting where the signal is
standardized, mandated, and institution-dependent. This distinction matters for
novelty. We do not claim to be the first paper to show that polarization can
shape financial disagreement. Instead, we argue that this paper is among the
first to show that political polarization changes how investors interpret
\emph{auditor-originated public signals}, and that the mechanism operates through
heterogeneous trust-based priors that are especially consequential when audit
signals are ambiguous.

The theoretical framework formalizes this mechanism. Investors observe the same
public audit signal but differ in their priors about the probability that the
signal reflects an underlying deterioration or improvement in reporting quality.
Those prior differences are linked to polarization through heterogeneous trust in
auditors and audit-regulatory institutions. In a Bayesian environment,
polarization generates a mean-preserving spread in priors, which increases the
dispersion of posterior beliefs after a common disclosure. Under standard
heterogeneous-belief trading logic \citep{harris1993, kandel1995, hong2007},
greater posterior dispersion implies more disagreement-driven trade and larger
absolute price responses. Importantly, the model yields a clean cross-sectional
prediction: polarization should matter most when the disclosure is harder to
interpret on its face, because posterior beliefs then depend more heavily on
priors and less on the signal itself.

This framing helps clarify what the paper contributes relative to adjacent
literatures. The audit-disclosure literature shows that auditor changes,
qualified opinions, going-concern opinions, and critical audit matters can
convey information to markets \citep{firth1978, dodd1984, francis2004,
defond2014}. The political-economy-of-auditing literature shows that political
forces can affect auditors, regulation, and reporting outcomes. And the newer
polarization-in-finance literature shows that politically divided environments
can generate disagreement and heterogeneous responses to common information.
What remains largely unaddressed is whether political polarization changes the
market interpretation of standardized audit disclosures. Item 4.01 events offer
an unusually clean setting to answer that question.

Empirically, we match auditor-change disclosures from Form 8-K to a measure of
political polarization based on election outcomes. The baseline design examines
whether absolute cumulative abnormal returns and abnormal trading volume around
auditor-change events are higher when the firm's local political environment is
more polarized. The main empirical claim is intentionally modest and focused:
polarization amplifies disagreement-type market responses to auditor changes.
The paper does not require signed average reactions to move in one universal
direction across all auditor changes, because auditor changes can be good news,
bad news, or ambiguous news depending on context. Instead, the paper predicts
that polarization increases the \emph{magnitude} of the reaction and the amount
of trade, with the strongest effects for disclosures whose informational content
is least clear ex ante.

This narrower thesis is better aligned with the strongest nearby literature.
Recent work links polarization to analyst disagreement \citep{hao2026} and to
broader financial disagreement and belief divergence \citep{kempf2024,
bordalo2023}. Those papers make it difficult to sustain a sweeping ``first
paper on polarization and disagreement'' claim. But they also make our audit
setting more compelling: if polarization affects how sophisticated market
participants process information generally, it is natural to ask whether the
same force changes how investors process auditor-originated disclosures that are
supposed to standardize market understanding. Likewise, prior auditor-change
studies establish that these events matter for capital markets \citep{johnson1990,
defond1998, landsman2009}, but they do not explain why market reactions to
similar disclosures vary systematically with the political environment.

The paper makes three contributions. First, it contributes to the auditing and
capital markets literature by showing that the informativeness of auditor-change
disclosures depends not only on event content, but also on the political context
in which investors receive the disclosure. Second, it contributes to the
polarization-in-finance literature by identifying a new mechanism ---
institutional-trust-based heterogeneity in the interpretation of an
external-auditor signal --- in a setting with standardized disclosure timing and
format. Third, it contributes to heterogeneous-beliefs research by bringing that
framework to an audit-specific empirical setting in which disagreement can be
observed through abnormal volume and absolute returns around mandated public
filings.

The paper also has regulatory implications. Much of modern audit regulation,
including expanded audit-reporting requirements, is motivated by the idea that
more standardized disclosure helps investors form more accurate and more uniform
beliefs. Our analysis suggests a limit to that logic. When investors differ in
their priors about institutional credibility, additional disclosure may not
resolve disagreement; in some cases it may even magnify disagreement by giving
investors more material to interpret through divergent priors. In that sense,
the effectiveness of disclosure-based regulation depends partly on the political
and institutional environment in which the disclosure is received.

The remainder of the paper proceeds as follows. Section~\ref{sec:lit} reviews
the related literature and positions the paper relative to adjacent work.
Section~\ref{sec:theory} develops the theoretical framework. Section~\ref{sec:hypotheses}
states the main testable hypotheses. Section~\ref{sec:design} describes the
empirical design and identification strategy. Section~\ref{sec:results}
provides the empirical tests and recommended analyses. Section~\ref{sec:conclusion}
concludes.

\section{Related Literature}
\label{sec:lit}

\subsection{Audit Signals in Capital Markets}

A long literature shows that auditor-originated disclosures matter for investors.
Early studies document market reactions to audit qualifications and related
reporting signals \citep{firth1978, dodd1984}. Subsequent work links audit
quality and auditor reputation to capital-market outcomes more broadly
\citep{francis2004, defond2014}. More recent research on expanded audit
reporting, including critical audit matters, studies whether additional
auditor-provided text improves market understanding or investor decisions
\citep{christensen2014, gutierrez2018}. This literature establishes that audit
signals can be value-relevant, but it largely treats the market's interpretation
of those signals as homogeneous or leaves cross-sectional variation in
interpretation unexplained.

Our paper shifts the question from whether auditor disclosures are informative on
average to whether their interpretation depends systematically on the political
environment. That distinction is important. If the same standardized disclosure
produces different levels of disagreement in different political environments,
then the market consequences of audit disclosure cannot be understood solely by
studying the signal in isolation.

\subsection{Auditor Changes as a Disclosure Setting}

Auditor changes are one of the cleanest settings in which to study the market
interpretation of audit-related information. Item 4.01 filings must be made
promptly, follow a relatively standardized structure, and often disclose whether
the registrant and former auditor disagreed on accounting or auditing matters.
Prior work documents that market reactions vary with the nature of the switch,
including whether the former auditor resigned or was dismissed, whether the
change appears to involve a movement in auditor quality, and whether
accompanying disclosures signal conflict or concern \citep{johnson1990,
defond1998, landsman2009}. We build on that literature but emphasize a different
margin: even after conditioning on event characteristics, investor reactions may
still vary because local political polarization changes how the event is
interpreted.

This framing helps sharpen the contribution. We are not simply re-documenting
that some auditor changes are worse than others. Rather, we ask why similarly
coded auditor-change events can generate different magnitudes of market response,
and we argue that one explanation lies in heterogeneous institutional trust in
polarized environments.

\subsection{Political Economy of Auditing and Regulation}

Another adjacent literature studies how political factors shape auditing,
reporting regulation, and enforcement. Existing work examines topics such as the
role of political connections, the effects of political turnover, and the impact
of politics on regulatory decision-making \citep{chaney2006, secshutdown2025,
ifrsideology2025}. This literature is highly relevant because it establishes that
auditing and financial reporting operate within a political environment rather
than outside of it. However, most of that work focuses on the \emph{supply side}
of audit quality, oversight, or enforcement. Our paper instead focuses on the
\emph{demand side}: how investors interpret an audit-related disclosure once it
is made.

This distinction is central for novelty. The paper does not argue that
polarization changes auditor behavior directly, although that channel may exist.
Instead, it isolates a different mechanism: polarization changes how investors
process a common auditor signal. That mechanism is especially plausible in audit
settings because the economic meaning of an auditor disclosure is mediated by
beliefs about the credibility of auditors, audit committees, the PCAOB, and the
broader regulatory architecture.

\subsection{Political Polarization in Finance}

A growing finance literature shows that polarization can affect beliefs, risk
perceptions, market expectations, and financial behavior \citep{kempf2024,
bordalo2023}. Particularly relevant is recent work showing that polarization is
associated with greater analyst disagreement \citep{hao2026}. That evidence is
important for this paper because it demonstrates that polarization can widen
belief dispersion among sophisticated financial intermediaries. At the same
time, it also means that our novelty claim must be carefully framed. We therefore
do not claim to be the first paper to connect polarization to financial
agreement or disagreement in general.

Instead, our contribution is narrower and more audit-specific. We examine whether
polarization changes the interpretation of \emph{auditor-originated public
signals} in a setting with standardized disclosure timing and mandated content.
This focus distinguishes the paper from existing polarization research centered
on analyst forecasts, trading behavior, media slant, or broad market beliefs.
It also allows us to contribute to the conference and journal audiences most
naturally interested in auditing, audit regulation, and capital markets.

\subsection{Heterogeneous Beliefs, Disagreement, and Common Signals}

Our theoretical mechanism draws on the heterogeneous-beliefs literature.
Differences of opinion can generate trading even without private information
\citep{harris1993}. When investors have heterogeneous priors, identical public
signals need not produce identical posteriors \citep{kandel1995}. More broadly,
research on disagreement and asset prices highlights how belief dispersion can
affect volume, volatility, and pricing outcomes \citep{hong2007}. Recent work on
disagreement over public-signal quality further suggests that common
information can widen rather than narrow disagreement when agents differ in how
credible or precise they perceive the signal to be \citep{divergingbeliefs2024,
publicsignalquality2024}.

This literature provides the conceptual bridge for our setting. Auditor-change
disclosures are common public signals. Political polarization creates a plausible
source of heterogeneity in priors about how those signals should be interpreted.
The combination implies a clear prediction: polarization should amplify
market-based measures of disagreement, especially when the disclosure is
ambiguous and therefore less capable of overwhelming prior differences.

\section{Theory}
\label{sec:theory}

We develop a simple framework that formalizes how political polarization can
change market reactions to auditor-change disclosures. The model is intentionally
parsimonious. Its purpose is not to explain the entire pricing of auditor-change
events, but to isolate the mechanism that is novel for this paper: greater
polarization increases heterogeneity in priors about the credibility and meaning
of an audit signal, thereby amplifying disagreement after a common disclosure.

\subsection{Setting}

Let $\theta \in \{0,1\}$ denote the latent state associated with an auditor-change
event, where $\theta=1$ represents a relatively favorable underlying state and
$\theta=0$ an unfavorable one. Investors observe a public signal $s$ from the
Item 4.01 disclosure. The signal may summarize facts such as whether the former
auditor resigned or was dismissed, whether disagreements were disclosed, and the
identity of the incoming auditor. To keep the model simple, suppose that after
parsing the filing, investors receive a coarse signal $s \in \{g,b\}$, where $g$
is more favorable than $b$.

Conditional on the latent state, signal precision is common knowledge:
\begin{align}
    \Pr(s=g\mid \theta=1) &= p_H, \\
    \Pr(s=g\mid \theta=0) &= p_L,
\end{align}
with $p_H > p_L$. The informativeness of the disclosure is governed by the gap
$p_H-p_L$. Smaller values correspond to more ambiguous disclosures.

A unit mass of investors observes the same public signal. Investors differ only
in their prior beliefs. Investor $i$ assigns prior probability $\pi_i \in (0,1)$
to the favorable state. We interpret variation in $\pi_i$ as arising from
heterogeneous trust in the institutional system surrounding auditor disclosures:
trust in auditors as monitors, trust in the firm's explanation of the change,
and trust in the regulatory institutions that make the disclosure credible.
Political polarization widens the cross-sectional dispersion of these priors.

\subsection{Polarization as a Source of Prior Dispersion}

To map the theory to the empirical setting, suppose investors belong to
politically differentiated groups indexed by $g \in \{A,B\}$. Group membership is
not important for its own sake; it captures the idea that more polarized
environments generate larger differences in institutional trust and therefore
larger differences in priors. Let the average prior remain fixed at $\bar{\pi}$,
but let polarization produce a mean-preserving spread in the distribution of
$\pi_i$.

This formulation is important for interpretation. The mechanism does not require
polarization to make investors uniformly more optimistic or more pessimistic
about auditor signals. It requires only that polarization makes investors less
likely to start from a common prior. That implication is exactly what makes
absolute returns and abnormal volume the natural primary outcomes.

\subsection{Belief Updating}

If the public signal is favorable ($s=g$), investor $i$'s posterior belief is
\begin{equation}
    \mu_i(g)
    = \Pr(\theta=1\mid s=g)
    = \frac{p_H\pi_i}{p_H\pi_i + p_L(1-\pi_i)}.
\end{equation}
If the signal is unfavorable ($s=b$), the posterior is
\begin{equation}
    \mu_i(b)
    = \Pr(\theta=1\mid s=b)
    = \frac{(1-p_H)\pi_i}{(1-p_H)\pi_i + (1-p_L)(1-\pi_i)}.
\end{equation}
In both cases, the posterior is monotone in the prior. Thus, a wider cross-sectional
dispersion of priors translates into a wider cross-sectional dispersion of
posteriors after the same public signal.

\begin{proposition}[Polarization increases posterior-belief dispersion]
\label{prop:dispersion}
Suppose the distribution of priors under high polarization is a mean-preserving
spread of the distribution of priors under low polarization. Then, for either
signal realization $s \in \{g,b\}$, the cross-sectional variance of posterior
beliefs is weakly higher under high polarization.
\end{proposition}

\begin{proof}
For either signal realization, the posterior is a monotone transformation of the
prior. A mean-preserving spread in priors therefore weakly increases the
cross-sectional spread of posteriors. \qed
\end{proof}

\subsection{Prices, Trading, and Absolute Reactions}

Let investor demand be increasing in the difference between the investor's
posterior valuation and the market price. In standard heterogeneous-belief
settings, greater dispersion in posterior beliefs increases the willingness of
investors to trade against one another, generating higher trading volume
\citep{harris1993, kandel1995}. Moreover, when the market-clearing price reflects
an aggregation of heterogeneous valuations, larger belief dispersion increases
the magnitude of the price response to a public signal for a given average shift
in beliefs, particularly when short-sale or risk-bearing constraints limit full
arbitrage.

These implications motivate our empirical focus on two disagreement-oriented
outcomes: abnormal trading volume and the absolute value of abnormal returns. In
contrast, the sign of the average return effect need not be common across all
auditor changes, because different event types can move average beliefs in
different directions.

\begin{proposition}[Polarization amplifies disagreement-oriented market outcomes]
\label{prop:market}
Holding constant the content of the auditor-change disclosure, higher political
polarization increases abnormal trading volume and the magnitude of the price
reaction to the disclosure.
\end{proposition}

\subsection{Signal Ambiguity}

The role of polarization should depend on how informative the disclosure is. If
$p_H-p_L$ is large, the signal is precise and investors' posteriors depend less
on their priors. If $p_H-p_L$ is small, the signal is ambiguous and priors matter
more.

\begin{proposition}[Ambiguity strengthens the polarization effect]
\label{prop:ambiguity}
The effect of prior dispersion on posterior-belief dispersion is stronger when
signal precision is lower. Consequently, the effect of political polarization on
abnormal volume and absolute returns is stronger for more ambiguous
auditor-change disclosures.
\end{proposition}

\begin{proof}
As signal precision falls, posterior beliefs place greater weight on priors.
Therefore, a given difference in priors translates into a larger difference in
posteriors. The result follows from the Bayes expressions above. \qed
\end{proof}

\subsection{Interpretation and Contribution}

The theory delivers three points that help strengthen the paper's contribution.
First, the mechanism is about the \emph{interpretation} of a standardized public
signal, not about private information or selective disclosure. Second, the main
predictions concern disagreement-type outcomes --- abnormal volume and absolute
returns --- which are the cleanest implications of heterogeneous priors in this
setting. Third, ambiguity is not a peripheral test; it is central to the theory
because it is the dimension along which trust-based priors should matter most.

This theoretical reframing also improves novelty. Rather than offering a general
``polarization matters'' story, the paper offers an audit-specific mechanism in
which polarization operates through heterogeneity in institutional trust and
therefore in the interpretation of auditor-originated disclosures.

\section{Hypothesis Development}
\label{sec:hypotheses}

The theory suggests that polarization should increase disagreement-based market
responses to auditor-change disclosures. We therefore focus the main hypotheses
on abnormal volume and absolute returns rather than on signed reactions.

\begin{hyp}
\label{h1}
Abnormal trading volume around auditor-change disclosures is higher in more
politically polarized environments.
\end{hyp}

\begin{hyp}
\label{h2}
The absolute cumulative abnormal return around an auditor-change disclosure is
higher in more politically polarized environments.
\end{hyp}

\begin{hyp}
\label{h3}
The positive relation between political polarization and disagreement-oriented
market reactions to auditor-change disclosures is stronger when the disclosure is
more ambiguous.
\end{hyp}

A signed-return prediction can still be explored as a secondary analysis within
subsamples where the direction of the signal is relatively clear (for example,
Big-4 downgrades, disclosed disagreements, or clearly favorable upgrades).
However, such tests should be framed as supplemental because the paper's main
mechanism concerns dispersion in beliefs, not a universal directional effect
across heterogeneous event types.

\section{Empirical Design}
\label{sec:design}

\subsection{Sample Construction}

\textbf{Auditor-change events.} We identify auditor changes using Form 8-K
filings disclosing a change in the registrant's certifying accountant under Item
4.01 for 2001--2023. We retain original 8-K filings that can be matched to CRSP
and Compustat and require sufficient return data to estimate abnormal returns.
The final sample consists of [PLACEHOLDER: N] auditor-change events for
[PLACEHOLDER: N] unique firms. We exclude regulated industries where return and
volume dynamics may differ systematically, including financial firms and
utilities, unless a later robustness test shows that including them does not
change the conclusions.

The paper should distinguish carefully among event subtypes: auditor resignation
versus dismissal, disclosed disagreement versus no disclosed disagreement,
movements toward or away from the Big 4, and changes whose quality direction is
unclear. These classifications are not merely descriptive controls; they are
central for testing the ambiguity mechanism.

\subsection{Dependent Variables}

\textbf{Absolute cumulative abnormal return ($|CAR|$).} We estimate abnormal
returns using a standard market model over the estimation window
$[-252,-46]$ trading days relative to the filing date. The main return outcome
is $|CAR(-1,+1)|$, which captures the magnitude of the market reaction while
remaining agnostic about sign.

\textbf{Abnormal trading volume ($AbVol$).} We measure abnormal trading volume as
the event-window mean of log turnover or log share volume relative to the firm's
pre-event baseline. Because theory predicts disagreement-driven trade, this
outcome is at least as important as the return-based specification and should be
treated as a co-equal primary dependent variable.

A useful extension would be to add one or both of the following: (i) abnormal
bid-ask spreads or intraday volatility as supplementary disagreement proxies, and
(ii) a signed-CAR analysis only within event subsamples for which prior literature
suggests a relatively clear direction.

\subsection{Political Polarization Measure}

The current draft uses an Esteban--Ray style polarization measure derived from
U.S. House election returns. That approach is reasonable, but the paper should
present it as a proxy for the extent to which the firm's local political
environment fosters divergence in institutional trust and interpretive priors.
The empirical discussion should therefore explain why headquarters-location
polarization is relevant: it captures the political environment most closely
associated with firm stakeholders, local media, employees, and a meaningful share
of the investor and attention base.

At the same time, the identification section should candidly acknowledge the
limits of a headquarters-based proxy for investor beliefs. To strengthen the
design, the paper should consider at least one of the following robustness
strategies:
\begin{enumerate}[leftmargin=*, itemsep=0.3em]
    \item replace state-level polarization with finer local measures where
    feasible;
    \item distinguish local polarization from national election-cycle effects
    using year fixed effects and additional controls;
    \item interact polarization with event ambiguity, which is less likely to be
    spuriously related to generic state characteristics; and
    \item perform placebo tests on non-audit 8-K events that are similar in
    disclosure timing but less dependent on trust in audit institutions.
\end{enumerate}

\subsection{Baseline Specification}

The core event-level specification is
\begin{equation}
\label{eq:main}
Y_i = \alpha + \beta_1 \, \textit{Polarization}_{s(i),t(i)}
+ \beta_2 \, \textit{Ambiguity}_i
+ \beta_3 \, \textit{Polarization}_{s(i),t(i)} \times \textit{Ambiguity}_i
+ \gamma'X_{i,t(i)-1} + \tau_{t(i)} + \delta_{j(i)} + \varepsilon_i,
\end{equation}
where $Y_i$ is either $|CAR|$ or $AbVol$, $\tau_{t(i)}$ are year fixed effects,
$\delta_{j(i)}$ are industry fixed effects, and $X_{i,t(i)-1}$ includes standard
firm controls such as size, leverage, profitability, market-to-book, prior stock
volatility, prior turnover, and indicators for financial distress or recent
reporting problems. Standard errors should be clustered at the firm level, with
a robustness check clustering by state or using two-way clustering when sample
structure permits.

The interaction term $\beta_3$ is important. In the current draft, ambiguity is
treated mainly through sample splits. That is useful descriptively, but the
paper will read as more disciplined if the ambiguity prediction is embedded
directly in the baseline specification and then reinforced through subsample
analyses.

\subsection{Identification and Threats to Inference}

The key empirical challenge is not whether polarization is measured with noise;
it is whether polarization is proxying for other local characteristics that also
affect market reactions to auditor changes. The draft should therefore discuss
and test the following concerns more explicitly.

\textbf{Local information environments.} More polarized states may have
different media environments, investor clienteles, or firm types. The paper
should include controls for local socioeconomic conditions and, where feasible,
firm-level information-environment proxies such as analyst coverage,
institutional ownership, or prior return volatility.

\textbf{Selection into auditor changes.} Polarization may affect the probability
that certain types of firms experience auditor changes. This concern does not
invalidate the event-study design, but it means the interpretation of cross-sectional
coefficients should be conditional on the occurrence of an auditor-change event.
A first-stage analysis of the incidence or type of auditor changes could be
helpful, even if only as a descriptive appendix.

\textbf{National political shocks.} Because polarization trends upward over time,
year fixed effects are necessary but may not be sufficient. The paper should show
that the results are not driven by a few election years or by the post-2016
period alone.

\textbf{Institutional-trust mechanism.} The strongest version of the paper would
include at least one proxy that makes the institutional-trust channel more
direct. Candidate moderators include measures of local trust in institutions,
firm-level governance quality, auditor reputation salience, or whether the filing
contains explicit disagreement language. Even if such evidence is imperfect, it
would substantially strengthen the mechanism.

\section{Results and Recommended Empirical Roadmap}
\label{sec:results}

Because the current draft is at a proposal stage, this section should be written
as a disciplined empirical roadmap rather than as a placeholder summary. The
paper will be stronger if the final version reports results in the following
sequence.

\subsection{Descriptive Evidence}

Begin by describing the sample composition and basic market responses around
Item 4.01 events. Report the distribution of event types, ambiguity measures, and
polarization. A simple figure plotting mean $|CAR|$ and abnormal volume by
polarization tercile would provide a transparent nonparametric first look.

\subsection{Baseline Tests of Hypotheses~\ref{h1} and \ref{h2}}

The first table should estimate Equation~\eqref{eq:main} without the ambiguity
interaction and then with it. The paper should treat volume and absolute returns
as parallel primary outcomes. The baseline interpretation should be straightforward:
if $\beta_1 > 0$, then auditor-change disclosures generate larger disagreement-type
market responses in more polarized environments.

To improve credibility, each table should move incrementally from a parsimonious
specification to one with rich firm controls, fixed effects, and alternative
clustering schemes. The text should emphasize economic magnitudes rather than
statistical significance alone.

\subsection{Ambiguity Tests of Hypothesis~\ref{h3}}

The ambiguity test should be central, not ancillary. The most compelling design
is an interaction between polarization and an event-level ambiguity indicator,
followed by targeted subsample analyses. Several ambiguity proxies are available:
\begin{enumerate}[leftmargin=*, itemsep=0.3em]
    \item no disclosed disagreement combined with no obvious quality-direction
    signal;
    \item same-tier auditor changes;
    \item sparse or boilerplate textual explanation in Item 4.01;
    \item cases in which the filing discloses a change but provides limited
    context about the reason.
\end{enumerate}

If the theory is correct, the interaction coefficient should be positive. This
result would sharply distinguish the paper from a generic story that polarized
states simply have more volatile stocks.

\subsection{Mechanism and Supplemental Tests}

The next set of analyses should show that the result is most pronounced when the
audit-institution channel is most salient. Candidate tests include:
\begin{itemize}[leftmargin=*, itemsep=0.3em]
    \item stronger effects for dismissals, disclosed disagreements, and Big-4
    downgrades or other auditor-quality-sensitive events;
    \item stronger effects for firms with weaker information environments, where
    investors have more room to rely on priors;
    \item weaker effects for comparatively clear and routine auditor changes;
    \item a signed-CAR analysis only within event categories whose directional
    content is well understood in prior research.
\end{itemize}

These tests would help transform the paper from a reduced-form association into
an audit-specific interpretation story.

\subsection{Robustness and Falsification}

A top-journal version of the paper should include several robustness checks.
First, vary the event window and the abnormal-return model. Second, use
alternative polarization measures where feasible. Third, re-estimate after
excluding influential states, election years, or the most politically salient
periods. Fourth, conduct placebo tests using nearby non-audit 8-K events.
Fifth, show that results are not driven entirely by generic return volatility or
turnover in polarized states.

A particularly valuable falsification test would compare auditor-change events to
other standardized 8-K disclosures that are economically important but less tied
to institutional trust in auditing. If polarization matters much more for
auditor-change filings than for those placebo disclosures, the paper's mechanism
becomes considerably more convincing.

\section{Conclusion}
\label{sec:conclusion}

This paper argues that political polarization changes how investors interpret
standardized auditor disclosures. The contribution is not a broad claim that
politics matters for markets in every setting, nor a generic replication of the
idea that polarization can widen disagreement. Instead, the paper identifies an
audit-specific mechanism: polarization increases heterogeneity in institutional
trust and therefore in prior beliefs about the meaning of auditor-originated
signals. In a setting where all investors observe the same mandated disclosure,
that mechanism predicts more disagreement-driven trade and larger absolute price
responses, especially when the signal is ambiguous.

This reframing improves both novelty and fit with the auditing literature. The
paper sits at the intersection of research on audit disclosures in capital
markets, political polarization in finance, and heterogeneous-belief asset
pricing, but its distinctive contribution is to show how these ideas interact in
auditor-change events. That positioning should resonate with a conference and
journal audience interested in auditing, capital markets, and the limits of
disclosure-based regulation.

In the next draft, the paper should prioritize four tasks: tightening the event
classification, making the ambiguity interaction central to the design,
strengthening the institutional-trust mechanism, and presenting the empirical
results in a way that foregrounds economic magnitude and falsification. If those
steps are executed well, the project can be framed as a credible and novel paper
for an auditing-and-capital-markets audience rather than as a broad political
finance paper with an audit application.

\begin{thebibliography}{99}

\bibitem[Bordalo et~al.(2023)]{bordalo2023}
Bordalo, P., Gennaioli, N., Ma, Y., and Shleifer, A. (2023).
\newblock Beliefs and economic behavior under polarization.
\newblock Working paper.

\bibitem[Chaney et~al.(2006)]{chaney2006}
Chaney, P., Faccio, M., and Parsley, D. (2006).
\newblock The quality of accounting information in politically connected firms.
\newblock Working paper / journal version as applicable.

\bibitem[Christensen et~al.(2014)]{christensen2014}
Christensen, B., Glover, S., and Wolfe, C. (2014).
\newblock Do critical audit matter paragraphs in the audit report change
nonprofessional investors' decision making?
\newblock Auditing-related working paper.

\bibitem[DeFond and Zhang(2014)]{defond2014}
DeFond, M. and Zhang, J. (2014).
\newblock A review of archival auditing research.
\newblock \emph{Journal of Accounting and Economics}, 58, 275--326.

\bibitem[DeFond et~al.(1998)]{defond1998}
DeFond, M., Wong, T., and Li, S. (1998).
\newblock The impact of improved auditor independence on audit market
concentration in China.
\newblock Placeholder entry; replace with the correct auditor-change citation used
in the final paper.

\bibitem[Dodd et~al.(1984)]{dodd1984}
Dodd, P., Dopuch, N., Holthausen, R., and Leftwich, R. (1984).
\newblock Qualified audit opinions and stock prices.
\newblock \emph{Journal of Accounting and Economics}, 6, 3--38.

\bibitem[Firth(1978)]{firth1978}
Firth, M. (1978).
\newblock Qualified audit reports: Their impact on investment decisions.
\newblock \emph{The Accounting Review}, 53, 642--650.

\bibitem[Francis(2004)]{francis2004}
Francis, J. (2004).
\newblock What do we know about audit quality?
\newblock \emph{British Accounting Review}, 36, 345--368.

\bibitem[Gutierrez et~al.(2018)]{gutierrez2018}
Gutierrez, E., Minutti-Meza, M., Tatum, K., and Vulcheva, M. (2018).
\newblock Consequences of adopting an expanded auditor's report in the United
Kingdom.
\newblock \emph{Review of Accounting Studies}, 23, 1543--1587.

\bibitem[Hao et~al.(2026)]{hao2026}
Hao, Q., Nain, A., and Xu, Y. (2026).
\newblock Political polarization and analyst disagreement.
\newblock Working paper.

\bibitem[Harris and Raviv(1993)]{harris1993}
Harris, M. and Raviv, A. (1993).
\newblock Differences of opinion make a horse race.
\newblock \emph{Review of Financial Studies}, 6, 473--506.

\bibitem[Hong and Stein(2007)]{hong2007}
Hong, H. and Stein, J. (2007).
\newblock Disagreement and the stock market.
\newblock \emph{Journal of Economic Perspectives}, 21, 109--128.

\bibitem[IFRS Ideology Paper(2025)]{ifrsideology2025}
Author(s) omitted. (2025).
\newblock Political ideology shapes reporting regulation.
\newblock Placeholder entry; replace with final bibliographic details.

\bibitem[Johnson and Lys(1990)]{johnson1990}
Johnson, W. and Lys, T. (1990).
\newblock The market for audit services: Evidence from voluntary auditor changes.
\newblock Placeholder entry; replace with final bibliographic details.

\bibitem[Kandel and Pearson(1995)]{kandel1995}
Kandel, E. and Pearson, N. (1995).
\newblock Differential interpretation of public signals and trade in speculative
markets.
\newblock \emph{Journal of Political Economy}, 103, 831--872.

\bibitem[Kempf et~al.(2024)]{kempf2024}
Kempf, E., McCartney, W., and others. (2024).
\newblock Political polarization and finance.
\newblock Working paper / review.

\bibitem[Landsman et~al.(2009)]{landsman2009}
Landsman, W., Nelson, K., and Rountree, B. (2009).
\newblock Auditor switches and market reactions.
\newblock Placeholder entry; replace with final bibliographic details.

\bibitem[Public Signal Quality Paper(2024)]{publicsignalquality2024}
Author(s) omitted. (2024).
\newblock Disagreement about public information quality and informational price
efficiency.
\newblock Placeholder entry; replace with final bibliographic details.

\bibitem[SEC Shutdown Paper(2025)]{secshutdown2025}
Author(s) omitted. (2025).
\newblock Economic consequences of political polarization: Evidence from an SEC
shutdown and its effect on insider trading.
\newblock Placeholder entry; replace with final bibliographic details.

\bibitem[Diverging Beliefs Paper(2024)]{divergingbeliefs2024}
Author(s) omitted. (2024).
\newblock Financial information and diverging beliefs.
\newblock Placeholder entry; replace with final bibliographic details.

\end{thebibliography}

\end{document}
